\input{../tex-inputs/latex-korrekturansicht-vorspann}

\section[Arthur Schnitzler an Marie Herzfeld, 19. 8. 1896]{L04328 Arthur Schnitzler an Marie Herzfeld, 19. 8. 1896}
\nopagebreak\mylabel{L04328v}
\rehead{ }\normalsize\beginnumbering\briefempfaengerindex{Herzfeld, Marie@\textsc{Herzfeld, Marie}!zzzSchnitzler, Arthur@\emph{von Arthur Schnitzler}!1896-08-191@{19. 8. 1896}|(be}
\toendnotes[C]{\smallbreak\pagebreak[2]}
\correspDesc{Versand  durch Arthur Schnitzler am 19. 8. 1896 in Skodsborg
\newline{}Übermittlung  am 19. 8. 1896 in Kopenhagen
\newline{}Zustellung  am 27. 8. 1896 in Grundlsee [Gemeinde]
\newline{}Erhalt  durch Marie Herzfeld am 27. 8. 1896 in Grundlsee [Gemeinde]}\toendnotes[C]{\smallbreak}
\Standort{keine Angabe, Privatbesitz, \emph{ohne Signatur}.}
\physDesc{Brief, 1 Blatt, 4 Seiten, Kuvert, 1418 Zeichen
\newline{}Handschrift: schwarze Tinte, deutsche Kurrent
\newline{}Versand: 1) Stempel: »\nobreak{}\oindex{Skodsborg@\textbf{Skodsborg}|pwk}\textcolor{gray}{Skodsborg}\nobreak{}«.   2) Stempel: »\nobreak{}\oindex{Kopenhagen@\textbf{Kopenhagen}|pwk}Kjobenhavn, 19. 8. 96, 5–6E\nobreak{}«.  3) Stempel: »\nobreak{}\oindex{Grundlsee [Gemeinde]@\textbf{Grundlsee [Gemeinde]}|pwk}Grundlsee, 27. 8. 96\nobreak{}«. 
\newline{}Ordnung: mit Bleistift Vermerk: »Zu 354« auf der Rückseite des Kuverts 
\newline{}Zusatz: 1) Zeitweise im Besitz von Gerhard Fichtner (1932–2012), vgl. die
                                 Erwähnung in seinem Aufsatz \emph{Rückblick auf »Leonardo«. Ein unbekannter Brief
                                       Sigmund Freuds vom 21. November 1932}. In: \emph{Luzifer–Amor}, H. 11, 1993,
                                    S. 135.  2) Versteigert bei \emph{Stargardt},
                                 Katalog 714 (Auktion 28. 4. 2026), Lot 163, um 750€.}\toendnotes[C]{\smallbreak}\pstart{}{\pb}\textsc{Dr. Arthur Schnitzler, \textcolor{pink}{Wien
                        IX}\oindex{IX., Alsergrund@\textbf{IX., Alsergrund}|pw}{}\ledrightnote{\textcolor{pink}{IX., Alsergrund}}.}{ }\textcolor{pink}{Frankgaſſe 1}\oindex{Wien@\textbf{Wien}!IX., Alsergrund@\textbf{IX., Alsergrund}!Frankgasse 1@\textbf{Frankgasse 1}|pw}{}\ledrightnote{\textcolor{pink}{Frankgasse 1}}.\pend{}{\bigskip}\pstart{}{\pb}\textsc{\textcolor{pink}{Steiermark – Austria}\oindex{Steiermark@\textbf{Steiermark}|pw}{}\ledrightnote{\textcolor{pink}{Steiermark}}}\pend{}\pstart{}\textsc{Fräulein Marie Herzfeld}\pend{}\pstart{}\textsc{\textcolor{pink}{Grundlsee}\oindex{Grundlsee [Gemeinde]@\textbf{Grundlsee [Gemeinde]}|pw}{}\ledrightnote{\textcolor{pink}{Grundlsee [Gemeinde]}}}\pend{}{\bigskip}\vspace{1em}
\pstart
           \raggedleft{}{\pb}\uline{d. Z. \textcolor{pink}{\textsc{Skodsborg}}\oindex{Skodsborg@\textbf{Skodsborg}|pw}{}\ledrightnote{\textcolor{pink}{Skodsborg}}},\pend
           
\pstart
           \raggedleft{}19. 8. 96.\pend
           
\pstart{}Verehrtes Fräulein,\pend\vspace{0.5em}
\pstart
           für Ihre ganz außerordentliche Liebenswürdgkeit danke ich Ihnen etwas verſpätet, denn
               Ihr \label{K_L04328-1v}\edtext{Brief}{\lemma{\textnormal{\emph{Brief}}}\Cendnote{\textnormal{Marie Herzfeld an Arthur Schnitzler, 7. 8. 1896.}}}\label{K_L04328-1} iſt mir erſt geſtern hier,
               in \textcolor{pink}{\textsc{Skodsborg}}\oindex{Skodsborg@\textbf{Skodsborg}|pw}{}\ledrightnote{\textcolor{pink}{Skodsborg}} zugeſtellt worden, wo ich ſeit ein paar Wochen bin. So konnte es ſich fügen, daſs ich \textcolor{blue}{\textsc{Georg Brandes}}\pwindex{Brandes, Georg 4.\,2.\,1842 Kopenhagen – 19.\,2.\,1927 ebd.@\textsc{Brandes, Georg} (4.\,2.\,1842 Kopenhagen – 19.\,2.\,1927 ebd.)|pw}{}\ledrightnote{\textcolor{blue}{Georg Brandes}} gerade an dem Tag, wo mir durch Sie {\pb}der Wortlaut ſeiner \textcolor{green}{Bemerkung}\pwindex{Brandes, Georg 4.\,2.\,1842 Kopenhagen – 19.\,2.\,1927 ebd.@\textsc{Brandes, Georg} (4.\,2.\,1842 Kopenhagen – 19.\,2.\,1927 ebd.)!To Forestillinger af Henrik IV@\strich\emph{To Forestillinger af Henrik IV}|pwv}{}\ledrightnote{{$\rightarrow$}\emph{\textcolor{green}{To Forestillinger af Henrik IV}}} über mich vermittelt
               wurde, perſönlich kennen lernte. Ich habe \label{K_L04328-2v}\edtext{geſtern bei ihm}{\lemma{\textnormal{\emph{geſtern bei ihm}}}\Cendnote{\textnormal{Vgl. A. S.: \emph{Tagebuch}, 18. 8. 1896.}}}\label{K_L04328-2} drei
               höchſt angenehme Stunden verbracht und viel amuſantes von ihm erfahren, das nie in
               eine Literaturgeſchichte ko{\geminationm}en wird, umſo literariſche
               Größen es ſich auch gehandelt hat.\pend
           
\pstart
           Was Sie, verehrtes Fräulein {\pb}mir von Ihren eignen Anſichten über die \textcolor{green}{L.}\pwindex{Schnitzler, Arthur 15. 5. 1862 Wien – 21. 10. 1931 ebd.@\textsc{Schnitzler, Arthur} (15. 5. 1862 Wien – 21. 10. 1931 ebd.)!Liebelei. Schauspiel in drei Akten@\strich\emph{Liebelei. Schauspiel in drei Akten}|pw}{}\ledrightnote{\textcolor{green}{Liebelei. Schauspiel in drei Akten}} mittheilen, iſt viel weniger, als ich gerade von
               Ihnen gern vernommen hätte. Es wird ſich vielleicht doch heuer in \textcolor{pink}{Wien}\oindex{Wien@\textbf{Wien}|pw}{}\ledrightnote{\textcolor{pink}{Wien}} Gelegenheit ergeben, 
               über das und auch über manches andre mit Ihnen zu reden; denn
               wie hoch ich die Klarheit ihres Urtheils und die ruhige Einfachheit ſchätze, mit der
               Sie es i{\geminationm}er {\pb}auszudrücken wiſſen, brauch ich Ihnen
               wohl nicht erſt zu verſichern. In der letzten Zeit habe ich ſo \label{K_L04328-3v}\edtext{wenig von Ihnen zu Gesicht beko{\geminationm}en}{\lemma{\textnormal{\emph{wenig … bekommen}}}\Cendnote{\textnormal{Im \emph{\textcolor{green}{Tagebuch}\pwindex{Schnitzler, Arthur 15. 5. 1862 Wien – 21. 10. 1931 ebd.@\textsc{Schnitzler, Arthur} (15. 5. 1862 Wien – 21. 10. 1931 ebd.)!Tagebuch@\strich\emph{Tagebuch}|pwk}} sind weder vor noch nach diesem Brief
                  Treffen ausgewiesen.}}}\label{K_L04328-3}; jetzt wird man hoffentlich \introOben{}bald\introOben{} wieder mehr von Ihnen leſen können.\pend
           
\pstart
           – Auch für ihre freundlichen Wünſche herzlichſten Dank; ſie erſchienen im rechten
               Moment!\pend
           
\pstart
           Mit verbindlichſten Grüßen{\\}Ihr Sie aufrichtig hochſchätzender{\\}\spacefill\mbox{Arthur Schnitzler}\pend
           \selectlanguage{ngerman}\endnumbering\briefempfaengerindex{Herzfeld, Marie@\textsc{Herzfeld, Marie}!zzzSchnitzler, Arthur@\emph{von Arthur Schnitzler}!1896-08-191@{19. 8. 1896}|)be}\mylabel{L04328h}
\begin{anhang}
\end{anhang}\newcommand{\dateiname}{L04328}\newcommand{\titel}{Arthur Schnitzler an Marie Herzfeld, 19. 8. 1896}\newcommand{\editorInnen}{Martin Anton Müller und Laura Untner}\input{../tex-inputs/latex-korrekturansicht-abspann}
